\section{实验设计与复现约束}
\subsection{历史实验的数据范围与证据分母}
工程基线固定于提交 \texttt{0af2de6a}，总计 1,881 个对象，
包括 1,769 个社区原始包与 112 个项目派生卡片。
本轮完整包审查只使用开发集 D1 的 200 个原始包，其中 107 个来自
\texttt{sickn33/agentic-awesome-skills}，93 个来自 \texttt{alirezarezvani/claude-skills}。
资源登记共 1,441 个文件，其中 334 个 Python 文件；339 个不支持文本分析或二进制资源保留为不透明内容。
现有留出材料混有项目卡片，来源连接组也很少，因此该批包审查不构成独立外部测试。

表\ref{tab:inventory}汇总不同实验的测量单位。
同一任务在多种库规模、检索模式、提示或重复下出现，不会增加独立任务数量。
独立人类质量金标准仍为零；本文新增的自动真值只覆盖明确定义的构造任务。

\begin{table*}[t]
\caption{前期历史实验及其推断范围。两项新成果另见后文；规模不是可互换的独立样本数。}
\label{tab:inventory}\centering\small
\begin{tabularx}{\textwidth}{p{.18\textwidth}p{.28\textwidth}X}
\toprule
实验 & 实际规模 & 主要测量与限制 \\
\midrule
完整包静态审查 & 200 个 D1 原始包，470 条命令观察 & 静态一致、冲突、未知；不测运行成功。\\
受控接口干预 & 35 个父包，381 个命令变体 & 154 个已知机械缺陷、227 个原始或无害对照；不是整体质量标签。\\
模型评价先导 & 18 个训练父包，三臂共 54 次请求 & 评分结构、引用定位、未知覆盖；没有独立语义人审。\\
来源补证诊断 & 389 条引用，另取两个许可资源 & 只核对真实来源窗口，保留模型摘录差异。\\
新产物实验 & 30 个输入，五臂两次重复，共 300 次 & 24 个 CSV 输入、6 个边界输入；一个目标技能、受限文本执行器。\\
网络选择回放 & 192 份捕获，六策略共 1,152 条记录 & 32 个既见任务；不测真实编排执行。\\
\bottomrule
\end{tabularx}
\end{table*}

\subsection{静态干预与真实模型评价}
从静态一致且标准库重建参照也接受的命令中，每个父包按固定哈希顺序至多取三条，
得到 80 条命令、35 个父包。
错误干预包括虚构选项、遗漏必填、非法枚举和数值类型错误；
无害对照包括原命令、空白变化和可等价使用的等号形式。
按照父包而不是命令变体重采样，避免重复样本制造精确区间。

模型评价使用 D1 训练 20 例中的 18 例，其余两例在响应产生之前因打包字节上限排除。
普通 rubric、支持证据优先和反证优先三臂使用同一包材料及锚点，首臂位置循环平衡。
每个响应必须包含八个不同维度，未知或不适用不能同时给出数值。
该先导是同模型提示诊断，不是 SkillNet 或 SkillCT 完整评价流程的复现。

\subsection{新构造任务与独立验收}
新的五臂实验在首次请求前冻结输入、答案、代码指纹、预算及比较方向。
任务按 12 个构造条件组织：普通记录、含逗号字段、字段内换行、空字段、字面 NA 字符串、
纯空白字段、完整重复、仅键重复、7/8 天临界、闰年日期、未来及缺失日期、混合引号。
每种条件各生成 8 行和 48 行两个输入，共 24 个范围内案例。
另有缺输入、不可读输入、XLSX、图片、无读取权限和请求删除六个元数据边界案例。
这些边界案例没有测试实际文件解码或操作系统权限执行。

任务要求输出行数、逐列空值数、完整重复的多余行数、超过七天的记录数、最大已知年龄，
并保留输入标识与空操作列表。
只有解析后的空字符串计为空值，NA、null 字面值和空格不计为空；
重复与过期按原始记录逐条计数，不自动删改。
答案由生成器的记录结构与年龄变量构造，再分别经标准库 CSV/日期解析与 pandas 读取核对。
24 个 CSV 案例三路一致；六个边界标签由明确的优先规则规定。
正确参考产物、逐字段错误、错误数值类型和多余字段共 390 项诊断全部通过。
两个读取实现仍共享任务语义，这不等于独立领域专家验证。

五臂分别为无 Skill、原始包、原始包加证据操作契约、无关技能上下文、仅操作契约。
目标为固定版本的 \texttt{data-quality-auditor}；原始上下文包含主文档及
\texttt{data\_profiler.py} 文本。
补充契约指出原脚本的 uniqueness=90、timeliness=85 是默认值而非测量，
并要求以任务定义为准逐项计算。它没有修改或执行社区脚本。
无关上下文取自天气资源生命周期、Solidity 安全和组织文化三个 D1 技能，并在整行字节上限内截取。
其长度接近原始上下文，但没有精确匹配实际 token 数。

\subsection{已完成历史模型实验的顺序、成本与统计}
两组真实模型实验均直连 DeepSeek 官方 API，使用 \texttt{deepseek-flash}，
关闭思考，温度为零，单次输出上限 2,048 token。
文档当时标注 DeepSeek-V4.1-Flash-1017；API 返回的是可变别名，不能作为不可变权重证明。
新实验并发为二；每个任务--重复组合下循环轮换五臂顺序，不依据结果重试或换模型。
实际完成全部 300 次，无基础设施失败及未运行项。

费用使用实际返回用量乘官方峰时、未缓存输入单价计算保守上界，
分别控制在 5 美元和 20 元人民币以内；两种币种的口径不是汇率换算。
缺失用量或费用状态未知时停止继续调度，并保留已在途结果。
账户余额差不作为逐次独立账单。

范围内结果按条件簇聚合，比较原始包相对无 Skill、证据契约相对原始包、
仅契约相对无 Skill、证据契约相对仅契约、无关上下文相对无 Skill。
对 12 个条件均值做 10,000 次重采样，给出条件百分位区间和配对胜负平。
这些条件来自目的性构造，不是独立随机抽取的真实任务族；区间不表征社区技能总体，
也不用于多重比较下的显著性或等效性认证。

网络回放则固定两个库规模、BM25/Hybrid/Fabric 三种模式和既见的 20/12 个任务。
Fabric 捕获只含图扩展，不含 LLM 重排。
选择器不接收参考标签，固定词/汉字双字组 TF-IDF 和名次推导的相关性。
全部 12 个库--任务--模式组别及六策略均保存，避免只报告有利单元。
